\section{Choosing the nonlinear and stationary data together}
\label{sec:common-funded-history}
The stationary construction, local expansion and scalar reset must use
one set of paths and one scale schedule. We now describe the simultaneous
choice and the way its identities are retained. The analytic inputs are
the path-attachment estimates, their extended-domain convergence bounds,
and the local nonlinear estimates proved in the preceding sections.
This section concerns their common application; a numerical choice of
large scales cannot replace an analytic estimate.

\subsection{The order of the choices}
Prescribe the finite derivative orders, spatial powers and support
exponents at every stage. They may tend to infinity. At a fixed stage
$j$, the adjacent target stage is either $j$ or its immediate predecessor;
thus every maximum below involves finitely many requirements.
First choose the free-kernel prices and block floors. The stationary
construction gives a positive upper width $\sigma_j$ and a lower path
allowance $A_j$. The startup expansion gives another lower path
allowance. Choose the profile's path allowance above their maximum,
and its width positive and at most $\sigma_j$. Keep its direct
coefficient below the purchased free-kernel bound. The profile
construction retains the prescribed orders and support exponents.
These choices precede the actual block lengths.

There are two levels of kernel price: the free construction's price
and the price required after replacing the terminal action by the
Wilson action. Choose the latter from the former now, at the fixed
profiles. Its validity is uniform over all later enlargements of the
block floors. Consequently the local nonlinear estimates can request
this final price without choosing it from a horizon that has not yet
been built.

Fix the startup and multiple-insertion requirements. Enlarge each
block floor by the maximum of the startup floor, its scalar comparison
floor and the multiple-insertion floor. Further enlarge it by the
floors for the nonlinear geometry and the auxiliary local comparisons.
The local targets have the same gauge chart, halo $500$, earned
exponent $b_i+3$, error exponent $-7/16$ and degree zero at target
stage $i$, where $b_i$ is the support exponent of stage $i$. These are
the exponents requested of the output. The target's own support
exponent is strictly larger than $b_i+3$ and its own rate is strictly
below $-7/16$; the two margins are used in
Section~\ref{sec:common-observable-purchases}. The relevant comparison
radius exceeds the halo by at
least five. These are the targets used to prove the interval return,
including the finite short-panel localization.

Choose any allowed block sequence above these floors. Only now select
the wage bounds and local price tables; their constants can depend on
the actual block lengths. An arbitrary larger wage family is still
allowed. Choose positive postal rates below any prescribed positive
upper rates, together with their forecast bounds. The rate agreement
needed by the scalar reset can be arranged explicitly. If $\rho_{ji}$
is the selected postal rate, take the post-step forcing rate $g_j$ no
larger than
\[
       \min\{\rho_{jj}/4,\rho_{j,j-1}/4\}.
\]
At stage zero the repeated predecessor is harmless. Set the reset
cost for the transition to $i$ equal to four times the local credit
constant at $i$. Thus
\begin{equation}\label{eq:common-reset-coverage}
             g_j\le\rho_{ji}/4,\qquad c_i\le q_{ji}/4
\end{equation}
at every permitted transition. These are positive quantities.
The post-step gap is supplied by the stationary estimate for this
actual block sequence; take the auxiliary multiplier equal to one
and the remaining post-step rate no larger than its stationary cap.

Next choose protection parameters below any prescribed positive
upper bounds and below one. The local expansion records can then be
selected with their smallness parameter at most $1/4$. The order is
significant: the interval return uses these records on every input
curve, so their selection precedes all parameter intervals and finite
periods.

Finally fix the insertion menus and auxiliary finite lists. Given any
previous entry delays, enlarge them successively to pay the following
already established estimates: the sector integrals and their active
Gaussian prefactors; the marked and unmarked local bounds; the two
forced absorptions; the small-field localization and parameter
comparison; and the scalar reset and startup prices. Impose both the
fixed local reserve $1/16$ and any requested positive smaller regular
reserve $\tau$, with forced reserve $1/8$. The support and path
conditions are identical in these two applications.

Here is a precise reason these successive enlargements are compatible.
Each estimate supplies a stage-indexed lower threshold $d^{(r)}_j$
and a conclusion valid whenever the entry coupling exceeds that
threshold. Replace the current threshold at stage $j$ by the finite
maximum of it and all the $d^{(r)}_j$ for that stage, and then by its
running maximum up to $j$. Shell lower bounds are treated the same
way. The conclusions persist at every later delay. Requirements on a
finite auxiliary list are imposed after that list is fixed, before the
final delay is selected. This proves the quantifier order
\[
 \text{profiles}\ ;\ \text{block lengths}\ ;\ \text{rates and gaps}\ ;\
 \text{protection}\ ;\ \text{menus and lists}\ ;\ \text{delays}\ ;\
 \text{scale realization}.
\]
The last realization is arbitrary among those paying these delays.
In particular no threshold is selected from a later finite-horizon
solution. Section~\ref{sec:common-scale-selection} constructs a
realization, with periods having arbitrarily late multiples of every
integer. This is the period divisibility needed for the five auxiliary
comparisons of the nonlinear return.

\subsection{Keeping the stationary and free data on that realization}
Let $a$ be one such realization. Build its free horizon $M^0$ with the
initial kernel prices, and its retained horizon $M^1$ with the next
prices. The retention identities keep their periodized covariance,
row and tower data in the precise combinations used by the stationary
construction. In particular the full-period lower bounds are supplied
for $M^1$, not for a separately selected quadratic form.
Apply the stationary construction to $M^1$ with the purchased gap and
shell bounds. It supplies simultaneously, for every cutoff $N$,
paths $R^1_{N,m}$, scalar corrections $\alpha_{N,m}$ and stationary
maps at every allowed period and forest. Its stronger version supplies
the extended-domain estimates used in
Section~\ref{sec:selected-history-limits}. The free convergence
of Section~\ref{sec:free-horizon-convergence} concerns these same
retained towers.

Replace only the terminal normalization, using the revised horizon
$M$ and paths $R_{N,m}$. The construction gives the exact identities
\begin{align*}
 R_{N,N}&=R^{\rm Wilson}_N,&
 R_{N,m}&=R^1_{N,m}\quad(m<N),\\
 \alpha^{\rm revised}_{N,m}&=\alpha_{N,m},&
 \rho^{\rm revised}_{N,m}&=\rho^1_{N,m}\quad(m+1<N).
\end{align*}
The entire tower agrees with that of $M^1$ when $m+1<N$, and its
transport matrix agrees already when $m<N$. The new horizon satisfies
the final nonlinear kernel price chosen before the block lengths.
At the terminal level its canonical direct coefficient is at most
one, which gives the required Wilson positivity bound. At lower levels
the retained all-period bounds apply.

For any fixed $m$, all the displayed identities hold on the tail
$N>m+1$. The limits and eventual larger-width bounds in
Section~\ref{sec:selected-history-limits} are therefore preserved:
remove finitely many terms from each sequence and use equality of the
remaining terms. This argument applies to each stationary derivative
as well as each path and tower coefficient. It also preserves the
extended-domain statement, because the one-step data coincide on this
tail. Thus the final nonlinear construction carries both the stationary
convergence data and the free-horizon convergence data. The terminal
normalization introduces no independently selected background.

\subsection{The same local return on every input interval}
On this final horizon the sector construction supplies its exact
Gaussian expression, including the active determinant, Haar factors,
sieve and coupling prefactor. The local price tables selected above
apply to this expression. The parameter comparison uses a single
reference dictionary and the identity identification of its elementary
labels, as in Section~\ref{sec:selected-nonlinear-offering}.
The auxiliary locality estimates are retained when the delays are
increased: they are stated beyond their entry thresholds, with all
block lengths, rates and local dictionaries fixed.

Consequently, for every cutoff, every adjacent pair of levels and
period, and every allowed continuous input curve on a compact interval,
that section supplies the exact return in one dictionary, with early
regular cost
\[
                 H=c_i\delta_E(\widehat\beta_m)
\]
and the forecast forced cost. Both the $1/16$ and the $\tau$ reserve
estimates refer to this same constructed horizon and local price data.
The selected tree and sector choice can depend on the interval and
period, as permitted by the scalar reset, but are fixed throughout
that interval.

We check the remaining recipe conditions rather than assuming that
they follow from the density identity. Write
\[
 \delta=\delta_E(\widehat\beta_m)/32,\qquad
 \gamma=\exp\{-\mathfrak k(\rho_{ji},\widehat\beta_{m+1},\beta)\}
\]
for the actual regular and forced prices after reorganization, where
$\mathfrak k$ is the purchased postal decay exponent and $\beta$ is
the incoming coupling. We verify the comparison with the recipe prices explicitly. Put
$x=\widehat\beta_m$, $y=x+c=\widehat\beta_{m+1}$, where
$c=s_G\log L\ge0$, and write $g(\beta)=\beta^{-1/2}$ and
$t(x)=g(x)P(x)$. Let $C\ge1$ be the exchange constant and let
$H_{\rm post}\ge1$, $D\in\mathbb N$, $\rho>0$ be the selected
postal constants. With $\ell_\theta(y)=P(y)^\theta$, define
\begin{align*}
 M(x,y)&=H_{\rm post}(1+B(x)+P(y)+2\ell_\theta(y))^D,\\
 R(x,y,\beta)&=M(x,y)\left[g(\beta)^3+g(\beta)^{-D}
 \left(e^{-\rho B(x)}+
 e^{-\rho P(y)^2/(1+B(x))^D}\right)\right],\\
 T(x,y,\beta)&=C\left[\left(\frac C L+\frac1{16C}\right)
                  \delta_E(y)+R(x,y,\beta)\right].
\end{align*}
All terms in $M$ and $R$ are nonnegative. The recipe's portion and
charges are precisely
\[
 \operatorname{portion}=T,\qquad
 \operatorname{charge}_E=T+M(x,y)t(x)^3,\qquad
 \operatorname{charge}_K=H_{\rm post}e^{-\mathfrak k},
 \quad \mathfrak k=\rho\left(\frac{t(y)}{g(\beta)}\right)^2.
\]
In particular $T\ge\delta_E(y)/16$. Since $c$ is fixed after $L$
is chosen, the elementary shifted-gap comparison gives, for all
sufficiently large $x$, $\delta_E(x)\le2\delta_E(y)$.
The reset constants of
Section~\ref{sec:physical-reset-reorganization} have the form
\[
 c_i=\frac1{64(1+B_i^{\rm reset})},\qquad B_i^{\rm reset}\ge0,
 \qquad B_i^{\rm reset}c_i\le\frac1{64},\qquad c_i\le\frac1{64}.
\]
Here $B_i^{\rm reset}$ is the coefficient of $H$ in the regular reset price;
write $v_i\ge0$ for the coefficient of $|\alpha|t(x)^3$ there.
It follows that
\[
 H=c_i\delta_E(x)\le\frac{\delta_E(y)}{32}\le T,
 \qquad
 \delta=\frac{\delta_E(x)}{32}\le\frac{\delta_E(y)}{16}
                                      \le T\le\operatorname{charge}_E.
\]
Moreover the uniform stage bound $|\alpha|\le A_j$ and the delayed
cubic estimate $64v_i A_jt(x)^3\le\delta_E(x)$ give
\[
 \operatorname{regularNeed}=B_i^{\rm reset}H+v_i|\alpha|t(x)^3
               \le\frac{\delta_E(x)}{64}
                   +\frac{\delta_E(x)}{64}=\delta.
\]
Thus the actual local cost is paid, as well as the recipe's larger
allowance. The soft reset price is bounded by a stage constant times
$(1+B(x))^4t(x)$; the late estimate for this expression in the reset
section bounds it by the forecast's mild reserve. Its remaining hard
reserve is exactly the purchased forecast comparison at this incoming
band. These two reserves are retained in the forced absorption of the
local return.

Since $H_{\rm post}\ge1$, $\gamma=e^{-\mathfrak k}\le
\operatorname{charge}_K$. This proves
\begin{equation}\label{eq:common-recipe-prices}
 H\le\operatorname{portion},\qquad
 \delta\le\operatorname{charge}_E,\qquad
 \gamma\le\operatorname{charge}_K.
\end{equation}
The incoming band lies in $[y/4,4y]$, so
$g(y)/2\le g(\beta)\le2g(y)$. Hence
$\mathfrak k\ge\rho P(y)^2/4\ge g_jP(y)^2$, using
\eqref{eq:common-reset-coverage}. Since $P(x+c)/P(x)\to1$ and $P(x)\to\infty$, eventually
$g_jP(x+c)^2\ge P(x)^{1/5}$. The coarse forced allowance is
$e^{-P(x)^{1/5}}$, so
\begin{equation}\label{eq:common-final-forcing}
 \gamma\le e^{-g_jP(y)^2}\le e^{-P(x)^{1/5}}.
\end{equation}
The coarse regular allowance dominates $\delta$ by its factor $32$.
Thus the output satisfies the costs of the next call as well as those
of the recipe batch.

For clarity, the corresponding implication on the actual arrays is
purely monotone. If $u_i$ and $v_i$ are their nonnegative regular and
forced majorants at costs $\delta,\gamma$, the same majorants work
at the larger charges in \eqref{eq:common-recipe-prices}; no supports,
functions or labels change. Section~\ref{sec:physical-reset-reorganization}
gives the exact normalized density identity for the same batch and
$t^2|\Delta|\le H\le\operatorname{portion}$. The permission predicate
used in this construction imposes no additional batch restriction.
Together these facts are exactly the batch relation of the chosen
recipe. Appending this batch to the incoming finite history supplies
the retained-history predicate at level $m$ and the next-step predicate.
This proves those predicates for the output that was actually
constructed.

The terminal startup construction, treated separately below, must supply one
continuous family on the full terminal band. Its regular and forced
bounds, common summable majorants and invariant jets survive the
fixed packing into the terminal dictionary. Its exact Wilson density
starts the retained history. This is the required initial-family assertion even
when $N=0$.

Once that startup assertion is proved, we can apply the actual scalar reset at every input interval
and then Section~\ref{sec:surviving-interval-shooting}. The result is a
finite history at every cutoff and period, with the prescribed
level-zero coupling, the exact transferred densities and the same
selected paths $R_{N,m}$. Its per-edge coupling error is the purchased
reset width, and its forced estimate includes
\eqref{eq:common-final-forcing}. No assertion of uniqueness of these
shooting histories is required. Their subsequent convergence and the
additional observable comparisons use the retained common data and
the estimates of their respective sections.
